\documentclass[11pt]{article}
\usepackage[margin=1in]{geometry}
\usepackage{amsmath, amssymb}
\usepackage{enumitem}
\usepackage{booktabs}
\usepackage{xcolor}
\usepackage{tcolorbox}

\title{\textbf{Poisson Distribution Crash Course}}
\author{}
\date{}

\begin{document}

\maketitle

\section{Definition \& Intuition}

The \textbf{Poisson distribution} models the \textbf{number of events} occurring in a fixed interval (time, space, volume, etc.) when events happen:
\begin{itemize}[noitemsep]
    \item Independently of one another
    \item At a constant average rate
\end{itemize}

\begin{tcolorbox}[colback=blue!5!white,colframe=blue!75!black,title=Key Insight]
Poisson answers: ``How many events will happen in this interval?''
\end{tcolorbox}

\section{PMF (Probability Mass Function)}

For $X \sim \text{Poisson}(\lambda)$:
\[
P(X = k) = \frac{\lambda^k e^{-\lambda}}{k!} \quad \text{for } k = 0, 1, 2, 3, \ldots
\]

where:
\begin{itemize}[noitemsep]
    \item $\lambda$ = average number of events in the interval
    \item $k$ = actual number of events we're calculating probability for
    \item $e \approx 2.71828$ (Euler's number)
\end{itemize}

\section{Mean \& Variance}

\begin{align*}
\mathbb{E}[X] &= \lambda \\
\text{Var}(X) &= \lambda \\
\text{SD}(X) &= \sqrt{\lambda}
\end{align*}

\begin{tcolorbox}[colback=red!5!white,colframe=red!75!black]
\textbf{Special Property:} Mean = Variance = $\lambda$
\end{tcolorbox}

\section{When to Use Poisson}

\subsection{Classic Examples}
\begin{itemize}
    \item Number of emails received per hour
    \item Number of customers arriving at a store per day
    \item Number of phone calls to a call center per minute
    \item Number of defects in a manufactured product
    \item Number of earthquakes in a region per year
    \item Number of mutations in a DNA sequence
\end{itemize}

\subsection{Requirements}
\begin{enumerate}
    \item \textbf{Independence:} Events occur independently
    \item \textbf{Constant rate:} Average rate $\lambda$ is constant throughout the interval
    \item \textbf{No simultaneity:} Two events cannot occur at exactly the same instant
\end{enumerate}

\section{Key Properties}

\subsection{Additivity}
If $X \sim \text{Poisson}(\lambda_1)$ and $Y \sim \text{Poisson}(\lambda_2)$ are independent, then:
\[
X + Y \sim \text{Poisson}(\lambda_1 + \lambda_2)
\]

\subsection{Approximation to Binomial}
When $n$ is large and $p$ is small (roughly $n \geq 20$ and $p \leq 0.05$):
\[
\text{Binomial}(n, p) \approx \text{Poisson}(np)
\]

\subsection{Scaling}
If $X \sim \text{Poisson}(\lambda)$ counts events in time $t$, then events in time $ct$ follow:
\[
\text{Poisson}(c\lambda)
\]

\section{Relationship to Exponential Distribution}

\begin{tcolorbox}[colback=green!5!white,colframe=green!75!black,title=Crucial Connection]
\textbf{Poisson} counts \textbf{how many} events in fixed time \\
\textbf{Exponential} measures \textbf{time between} consecutive events
\end{tcolorbox}

\vspace{0.3cm}

If events arrive as a Poisson process with rate $\lambda$ per unit time:
\begin{itemize}
    \item Number of events in time $t$: $\sim \text{Poisson}(\lambda t)$
    \item Time between consecutive events: $\sim \text{Exponential}(\text{mean } = 1/\lambda)$
\end{itemize}

\subsection{Exponential Distribution Recap}
For $T \sim \text{Exponential}(\text{mean } = \mu)$:
\begin{align*}
\text{PDF: } f(t) &= \frac{1}{\mu} e^{-t/\mu} \quad \text{for } t \geq 0 \\
\text{CDF: } F(t) &= 1 - e^{-t/\mu} \\
\text{Survival: } S(t) &= e^{-t/\mu}
\end{align*}

Alternative parameterization with rate $\lambda = 1/\mu$:
\begin{align*}
\text{PDF: } f(t) &= \lambda e^{-\lambda t} \\
\text{Survival: } S(t) &= e^{-\lambda t}
\end{align*}

\section{Worked Examples}

\subsection{Example 1: Basic Calculation}
A call center receives an average of $\lambda = 3$ calls per minute.

\textbf{Q:} What's the probability of exactly 2 calls in the next minute?

\textbf{Solution:}
\begin{align*}
P(X = 2) &= \frac{3^2 e^{-3}}{2!} \\
&= \frac{9 \times 0.0498}{2} \\
&= \frac{0.4482}{2} \\
&= 0.224
\end{align*}

\subsection{Example 2: Zero Events}
With $\lambda = 3$ calls per minute, what's the probability of \textit{no} calls?

\textbf{Solution:}
\begin{align*}
P(X = 0) &= \frac{3^0 e^{-3}}{0!} \\
&= e^{-3} \\
&\approx 0.0498
\end{align*}

\subsection{Example 3: At Least One Event}
What's the probability of at least 1 call?

\textbf{Solution:}
\begin{align*}
P(X \geq 1) &= 1 - P(X = 0) \\
&= 1 - e^{-3} \\
&\approx 0.950
\end{align*}

\subsection{Example 4: Time Between Events}
If calls arrive as Poisson($\lambda = 3$ per minute), what's the average time between calls?

\textbf{Solution:}
Time between calls $\sim$ Exponential(mean $= 1/\lambda = 1/3$ minute $= 20$ seconds)

\textbf{Bonus:} What's the probability the next call arrives within 30 seconds (0.5 minutes)?
\begin{align*}
P(T \leq 0.5) &= 1 - e^{-\lambda t} \\
&= 1 - e^{-3 \times 0.5} \\
&= 1 - e^{-1.5} \\
&\approx 0.777
\end{align*}

\subsection{Example 5: Scaling}
If $\lambda = 3$ calls per minute, how many calls do we expect in 5 minutes?

\textbf{Solution:}
Number of calls in 5 minutes $\sim$ Poisson($5 \times 3 = 15$)

Expected number: $\mathbb{E}[X] = 15$ calls

\section{Common Pitfalls}

\begin{enumerate}
    \item \textbf{Don't forget $k!$ in denominator} - It's easy to drop when calculating by hand
    \item \textbf{Rate vs. mean confusion} - In exponential: if Poisson rate is $\lambda$, exponential mean is $1/\lambda$
    \item \textbf{Scaling the rate} - If $\lambda = 2$ per hour, then for 3 hours use $\lambda = 6$
    \item \textbf{Independence assumption} - Poisson requires events to be independent (e.g., doesn't work well for bus arrivals that run on a schedule)
\end{enumerate}

\section{Quick Reference Table}

\begin{table}[h]
\centering
\begin{tabular}{@{}ll@{}}
\toprule
\textbf{Property} & \textbf{Formula/Value} \\
\midrule
PMF & $P(X = k) = \frac{\lambda^k e^{-\lambda}}{k!}$ \\
Mean & $\lambda$ \\
Variance & $\lambda$ \\
Standard Deviation & $\sqrt{\lambda}$ \\
$P(X = 0)$ & $e^{-\lambda}$ \\
$P(X \geq 1)$ & $1 - e^{-\lambda}$ \\
Related Exponential Mean & $1/\lambda$ \\
\bottomrule
\end{tabular}
\end{table}

\end{document}
